\documentclass[12pt]{article}
\usepackage[margin=2.5cm]{geometry}
\usepackage{sintaxismark}
\usepackage{sintaxismark-grid}

\title{\texttt{sintaxismark} v0.7.0 -- grid mode demo}
\author{}
\date{}

\begin{document}
\maketitle

\section*{1. New experimental grid syntax}

The new mode uses the same labels inside a dedicated environment. The commands do not draw immediately; the environment extracts tokens, infers ranges, computes line breaks, and then draws the analysis.

\begin{sintaxisgrid}[width=12cm]
\SA{
  \SES{Juan}
  \PVS{
    vio
    \OD*{
      \MD{un}
      \N{tigre}
      \NN{orgulloso}
      \MI*{
        de
        \COMP*{
          \MD{sus}
          \N{garras}
        }
      }
    }
  }
}
\end{sintaxisgrid}

\section*{2. Long example with automatic line breaking}

\begin{sintaxisgrid}[width=13cm]
\SA{
  \PVS{
    Te propongo
    \OD*{
      \MD{una}
      \N{consigna}
      \MI*{
        de
        \COMP*{
          \N{comprobacion}
          \MI*{de \N{lectura}}
        }
      }
      \MI*{
        sobre
        \COMP*{
          \MD{la}
          \N{clasificacion}
          \NN{general}
          \MI*{de \N{Bosque}}
        }
      }
    }
    con ejemplos propios o tomados del texto
  }
}
\end{sintaxisgrid}

\section*{3. Starred and unstarred behavior}

\begin{sintaxisgrid}[width=12cm]
\SA{
  \PVS{
    analizo
    \OD{
      \MD{esta}
      \N{oracion}
      \MI*{con \N{cuidado}}
    }
  }
}
\end{sintaxisgrid}

\bigskip
\noindent\textbf{Legend.} Inside \texttt{sintaxisgrid}, upper labels such as \texttt{SES} and \texttt{PVS} become upper spans; starred lower labels such as \texttt{OD*}, \texttt{MI*}, and \texttt{COMP*} become lower open boxes; unstarred \texttt{OD} and \texttt{OI} become underline+label; labels such as \texttt{MD}, \texttt{N}, \texttt{NN}, and \texttt{NV} become simple tags.

\end{document}
